% cover.tex: the cover page of the ECN 593 final report.
% Every result number is a macro from numbers.tex (scripts/write_numbers.py).
% The diagram uses data coordinates. The axis ranges and label offsets are
% layout: check the cover by eye after a rerun that moves a value.
\begin{titlepage}
\thispagestyle{empty}
\setlength{\parskip}{0pt}
\raggedright

{\footnotesize\color{line}%
ECN 593 \textperiodcentered\ APPLIED PROJECT \textperiodcentered\ FINAL REPORT\hfill ARIZONA STATE UNIVERSITY\par}
\vspace{0.5em}
{\color{rule}\rule{\linewidth}{1.2pt}}
\vspace{1.2em}

{\fontsize{24}{29}\selectfont\bfseries\color{ink}%
Not-ruled-out images\\[0.15em]of policy contrasts\par}
\vspace{0.9em}
{\normalsize\color{line}%
A village agent-based model, calibrated by history matching,\\
and two analysts (a field experimenter and a calibrator)\\
graded against what the targets cannot rule out\par}

\vspace{1.2em}

% ---- the paper in one picture: the min-I point (the best-fitting NROY point)
{\normalsize\bfseries\color{ink} The paper in one picture\par}
\vspace{0.7em}
\begin{tikzpicture}[
    every node/.style={font=\footnotesize, text=ink},
    band/.style={fill=wash, draw=line, line width=0.3pt},
    foil/.style={line width=2.4pt, draw=ink},
    truth/.style={font=\normalsize},
    tick/.style={line width=0.4pt, draw=line}]
  % Row 1: theta_2 on [0, 0.08] mapped to 16.4 cm (x = 205 cm per unit)
  \begin{scope}[x=205cm, y=1cm]
    \node[anchor=south west] at (0, 0.78) {\bfseries Foil 1\normalfont\ \ ignores interference \textperiodcentered\ policy effect};
    \fill[band] (\ThetaHmLo, 0.18) rectangle (\ThetaHmHi, 0.74);
    \node[anchor=north east, text=line] at (\ThetaHmHi, 0.72) {image $\Theta_{\mathrm{hm}}$};
    \draw[foil] (\MinIFoilOneLo, 0.33) -- (\MinIFoilOneHi, 0.33);
    \node[truth] (star1) at (\MinIThetaTwo, 0.56) {$\star$};
    \node[anchor=west] at (star1.east) {truth here: \MinIThetaTwoShort};
    \draw[tick] (0, 0) -- (0.08, 0);
    \foreach \v in {0, 0.02, 0.04, 0.06, 0.08} {
      \draw[tick] (\v, 0) -- (\v, -0.07);
      \node[below, text=line] at (\v, -0.05) {\v};}
  \end{scope}
  % Row 2: delta_spill on [-0.004, 0.028] mapped to 16.4 cm (x = 512.5 cm per unit)
  \begin{scope}[x=512.5cm, y=1cm, xshift=2.05cm, yshift=-2.05cm]
    \node[anchor=south west] at (-0.004, 0.78) {\bfseries Foil 2\normalfont\ \ models exposure \textperiodcentered\ spillover};
    \fill[band] (\ThetaExpLo, 0.18) rectangle (\ThetaExpHi, 0.74);
    \node[anchor=north east, text=line] at (\ThetaExpHi, 0.72) {image $\Theta_{\mathrm{exp}}$};
    \draw[foil] (\MinIFoilTwoLo, 0.33) -- (\MinIFoilTwoHi, 0.33);
    \node[truth] (star2) at (\MinISpill, 0.56) {$\star$};
    \node[anchor=west] at (star2.east) {truth here: \MinISpillShort};
    \draw[tick] (-0.004, 0) -- (0.028, 0);
    \foreach \v in {0, 0.01, 0.02} {
      \draw[tick] (\v, 0) -- (\v, -0.07);
      \node[below, text=line] at (\v, -0.05) {\v};}
  \end{scope}
\end{tikzpicture}

\vspace{0.4em}
{\footnotesize\color{line}%
The diagram shows the min-I point, the calibrator's best-fitting NROY point. The two
rows use different scales. Bands: the image of the effect over \ThetaSample{} sampled NROY points. Black
lines: the analyst's nominal \ConfidencePct\% interval from one panel. Stars: the
simulator's true value at this point. Foil~1 is tight, sits inside the band,
and misses the truth. Foil~2 contains the truth, yet it also sits inside its
band: narrower than what the targets leave open.\par}

\vspace{0.9em}

% ---- question and answer
\noindent
\begin{minipage}[t]{0.47\linewidth}\raggedright
{\normalsize\bfseries\color{ink} The question}\par\vspace{0.4em}
{\small Many parameter points fit the same \NMoments{} targets (the status quo plus
one randomized pilot), and they imply different effects of seeding the
hubs. What must a calibrated simulator report? How far can an interval
computed at one point be trusted?\par}
\end{minipage}\hfill
\begin{minipage}[t]{0.47\linewidth}\raggedright
{\normalsize\bfseries\color{ink} The answer}\par\vspace{0.4em}
{\small In the one panel at each point, Foil~1 misses the truth at all
three points, the real village included, and is F-OVERCONF at all three.
Foil~2 contains the truth at all three, and is F-OVERCONF at both
calibrated points; at the real village it has no verdict. Report the image,
not one point.\par}
\end{minipage}

\vspace{0.9em}

% ---- in numbers
{\normalsize\bfseries\color{ink} In numbers\par}
\vspace{0.4em}
{\footnotesize
\begin{tabularx}{\linewidth}{@{}X r@{}}
\toprule
\textcolor{line}{Points evaluated in three history-matching waves / points kept (NROY)} & \PointsEvaluated{} / \NNroy{} \\
\textcolor{line}{Image of the policy effect over the sampled NROY points} & [\ThetaHmLoShort, \ThetaHmHiShort] \\
\textcolor{line}{Policy effect at the planted point c*} & \ThetaTwoStarShort{} \\
\textcolor{line}{Foil 1 estimate below the truth, at $c^\star$ / min-I / median-I} & \StarPtFoilOneBelowPct\% / \MinIFoilOneBelowPct\% / \MedIFoilOneBelowPct\% \\
\textcolor{line}{Foil 2 interval coverage when the exposure map holds exactly} & \GateCoveragePct\% \\
\bottomrule
\end{tabularx}\par}

\vfill

{\color{rule}\rule{\linewidth}{0.4pt}}\par\vspace{0.6em}
\noindent
\begin{minipage}[t]{0.44\linewidth}\raggedright
{\small\bfseries\color{ink} Augusto Gonz\'alez-Bonorino}\par\vspace{0.15em}
{\footnotesize\color{line} Arizona State University\\
Applied project for ECN 593\par}
\end{minipage}\hfill
\begin{minipage}[t]{0.54\linewidth}\raggedleft
{\footnotesize\color{line} Draft for review, October 2026\\
\texttt{github.com/Bonorinoa/ecn593-applied-project}\\
Deterministic \textperiodcentered\ 55 tests \textperiodcentered\ one command replicates it\par}
\end{minipage}
\end{titlepage}
